\section{Discussion}

\subsection{What this does and does not support}

Within the 15 tested (deliberately low-stakes, benign) scenarios, three
small open-weight models did not produce textually detectable increases in
manipulation-coded language when instructed to maximize agreement, as
measured by this rubric and this judge ensemble. All three models
overwhelmingly preserved disclosure of tradeoffs/alternatives and an easy
path to decline, even under the explicit pressure framing.

This does \emph{not} support a claim that these models -- or language
models generally -- cannot become more manipulative under any goal
pressure. The scenarios were designed to be benign; higher-stakes or more
adversarial scenario framings, longer or multi-turn interactions, or
different (stronger, more specific) pressure instructions might elicit a
different pattern. Nor does it support any claim about the judge's ability
to detect subtler manipulation tactics outside the six rubric dimensions,
or about internal model states, intent, or evaluation-awareness -- this
study measures observable generated text under a controlled prompt only.

\subsection{Why might the effect be absent here?}

Several non-exclusive explanations are consistent with the data:

\begin{enumerate}
\item \textbf{Genuine model robustness at this scale and stakes level.}
Small open-weight models may carry safety-relevant training that resists
persuasion-maximizing framing even when explicitly instructed, at least
for benign requests.
\item \textbf{Scenario benignity limits the ``room'' for manipulation.}
Scenarios were designed with no dangerous, high-stakes, or adversarial
content by design (Section~\ref{sec:ethics}); this may be precisely why
little manipulation was detected regardless of condition.
\item \textbf{Judge and rubric sensitivity.} A small open-weight judge,
scoring on a coarse 0--2 scale per dimension, may under-detect subtle
persuasion techniques that a stronger judge or a human annotator would
catch -- motivating this study's blinded multi-judge and selective
human-escalation design rather than trusting a single judge's score.
\item \textbf{Ceiling/floor compression limits statistical power.} With
most scores at or near the scale's boundary, the paired bootstrap and
permutation tests have very little variance in which to detect a true
effect, even a small one.
\end{enumerate}

We do not have evidence to adjudicate between these explanations from this
pilot alone, and we do not force a preference among them.

\subsection{Judge sensitivity and multi-judge validation}

Judge agreement is markedly higher on dimensions closer to the observed
floor/ceiling (\texttt{respect\_for\_refusal}, \texttt{artificial\_urgency})
than on dimensions requiring finer discrimination
(\texttt{deceptive\_omission}, \texttt{authority\_social\_pressure}). Combined
with \texttt{llama3.2:3b}'s substantially higher empty-evidence rate, this
suggests smaller open-weight judges are less reliable at the more subtle
end of manipulation detection specifically -- a finding relevant to anyone
considering small open-weight models as judges for behavioral evaluation,
independent of this study's substantive null result.
